\documentclass[UTF8]{ctexart}
\usepackage{lmodern} % 使用准公文体
% 导入所需的宏包
\usepackage{array,booktabs} % 表格增强
\usepackage{amsmath} % 数学公式
\usepackage{graphicx} % 图片插入
\usepackage{cite} % 引用格式
\usepackage{listings} % 代码插入
\usepackage{hyperref}
\usepackage{xcolor} % 颜色
\usepackage{tikz} % 绘图
\usepackage{fontspec} % 字体设置
\usepackage{adjustbox} % 调整盒子大小
\usepackage{subcaption} % 子图
\usepackage{url} % 超链接
\usepackage{multirow}
% \setCJKmainfont{SimSun}[AutoFakeBold=3.0]

% 设置页面布局
\usepackage[top=1.5cm, bottom=1.5cm, left=0.5cm, right=0.5cm]{geometry}
\definecolor{darkred}{HTML}{8B0000}
\definecolor{darkblue}{HTML}{0044FF}

\usepackage{mdframed}
\usepackage{titlesec}
\usepackage{indentfirst} % 用于调整段落格式
\titleformat{\paragraph}
{\normalfont\small\bfseries} % 字体大小为\normalsize，加粗
{\theparagraph}{1em}{} % 编号与标题之间的间距为1em
\titlespacing*{\paragraph}
{16pt}{3.25ex plus 1ex minus .2ex}{1ex plus .2ex} % 段前间距和段后间距

\titleformat{\subparagraph}
{\normalfont\small\bfseries} % 字体大小为\normalsize，加粗
{\theparagraph}{1em}{} % 编号与标题之间的间距为1em
\titlespacing*{\subparagraph}
{16pt}{3.25ex plus 1ex minus .2ex}{1ex plus .2ex} % 段前间距和段后间距

\newenvironment{boxedenum_green}{
  \begin{mdframed}[
    backgroundcolor=green!5,
    linecolor=black,
    outerlinewidth=2pt,
    roundcorner=10pt,
    innertopmargin=10pt,
    innerbottommargin=10pt,
    innerleftmargin=-10pt,
    % leftmargin=1.5em,
    rightmargin=1.2em,
  ]
  \begin{enumerate}
}{
  \end{enumerate}
  \end{mdframed}
}

\newenvironment{boxedenum_gray}{
  \begin{mdframed}[
    backgroundcolor=gray!10,
    linecolor=black,
    outerlinewidth=2pt,
    roundcorner=10pt,
    innertopmargin=10pt,
    innerbottommargin=10pt,
    innerleftmargin=0pt, % 取消框内侧的左侧缩进
    leftmargin=1.5em, % 左侧外部缩进
    rightmargin=1.2em, % 右侧外部缩进
  ]
  \begin{enumerate}
}{
  \end{enumerate}
  \end{mdframed}
}

\setcounter{topnumber}{10} % 每页顶部允许的浮动体数量
\setcounter{bottomnumber}{10} % 每页底部允许的浮动体数量
\setcounter{totalnumber}{20} % 每页允许的总浮动体数量
\renewcommand{\topfraction}{1.0} % 页面顶部允许的最大浮动体比例
\renewcommand{\bottomfraction}{1.0} % 页面底部允许的最大浮动体比例
\renewcommand{\textfraction}{0.0} % 页面上非浮动体部分的最小比例
\renewcommand{\floatpagefraction}{1.0} % 浮动体页面的最小比例


%================== 进度条（无 collectcell）==================%
\newlength{\progwidth}\setlength{\progwidth}{6cm}
\newlength{\progheight}\setlength{\progheight}{8pt}

% 安全版颜色计算
\newcommand{\progfcolor}[1]{%
  \pgfmathsetmacro{\tmpA}{#1*2}%
  \ifdim\tmpA pt>1pt\relax
    \pgfmathsetmacro{\tmpB}{255-(#1-0.5)*510}%
    \definecolor{progtmp}{RGB}{\tmpB,255,0}%
  \else
    \pgfmathsetmacro{\tmpC}{#1*510}%
    \definecolor{progtmp}{RGB}{255,\tmpC,0}%
  \fi
}

% 进度条 + 百分比
\newcommand{\progressbar}[1]{%
  \progfcolor{#1}%
  \begin{tikzpicture}[baseline=-0.5ex]
    \fill[gray!30] (0,0) rectangle (\progwidth,\progheight);
    \fill[progtmp]   (0,0) rectangle (#1\progwidth,\progheight);
    \node[anchor=west,inner xsep=2pt,font=\sffamily\small] at (0,\progheight/2)
         {\pgfmathparse{round(#1*100)}\pgfmathresult\%};
  \end{tikzpicture}%
}
%================== 表格列类型 ==================%
\newcolumntype{P}{>{\raggedright\arraybackslash}p{7cm}} % 进度条列宽自己定


% 设定文档标题、作者和日期
\fontsize{10}{12}\selectfont
\title{\textbf{两周小结与计划表 \\ {\fontsize{12pt}{4pt}\selectfont (Two-Week Research Summary and Plan)}}}
\author{
  \makebox[0.5\textwidth][l]{姓名(Name)：白春辉}
  \makebox[0.4\textwidth][r]{日期(Date)：\large\today}
%   \makebox[0.4\textwidth][r]{日期(Date)：\large 2024年12月29日}
}
\makeatletter
\renewcommand{\@date}{} % 将 \@date 定义为空，从而不显示日期
\makeatother
\begin{document}

\maketitle % 输出标题、作者和日期
\vspace{-3.5em}
\centerline{\rule{\linewidth}{1pt}}
\vspace{-1em}
\section*{\textbf{过去两周的研究小结 \\ {\fontsize{12pt}{4pt}\selectfont (Research work done in the last two weeks)}}}

% \begin{table}[htbp]
% \centering
% % \caption{两周任务进度}
% \begin{tabular}{@{}p{4cm}ll@{}}
% \toprule
% \textbf{工作任务} & \textbf{进度情况} & \textbf{未完成情况} \\
% \midrule
% 论文初稿撰写 & \progressbar{1.0} &  \\
% 各老师审阅\&修改 & \progressbar{1.0} &  \\
% \bottomrule
% \end{tabular}
% \end{table}

\section*{基于状态距离矩阵诱导表征的 Trajectory Transformer 技术方案}

\subsection*{1. 目的与动机 (Motivation)}
在强化学习 (RL) 环境中，传统的价值函数估计或策略梯度方法往往侧重于即时奖励优化。本方案将 \textbf{RL 任务考虑为序列建模问题 (Sequence Modeling Problem)}，旨在通过序列预测任务实现面向长时序目标的决策。
\begin{itemize}
    \item \textbf{建模范式}：借鉴Decision Transformer (DT) 的思想，将智能体与环境交互产生的轨迹视为一种可预测的状态序列，通过自回归建模预测状态演化。
    \item \textbf{核心挑战}：强化学习的状态空间通常呈现出高维、连续且规模庞大的特征，其核心瓶颈在于状态内部存在巨大的语义鸿沟：由于缺乏天然的词表约束，模型极难界定两个状态的本质相似性，往往会出现特征空间距离相近但物理语义截然不同的现象。此外，实际问题的状态轨迹具有极强的局部连续性，而 Transformer 固有的全局注意力机制在训练初期往往难以捕捉这种时空连续性约束，从而生成与实际不符的非法轨迹。
    \item \textbf{解决方案}：基于度量空间同构化的嵌入层初始化。本方案提出一种基于多维尺度变换（MDS）的嵌入空间预构建方法，本质是一种领域知识增强的嵌入层初始化。其核心逻辑在于：
    \begin{itemize}
        \item 拓扑度量化：首先定义一种问题特定的状态距离度量，将环境的物理约束或逻辑代价显式化；
        \item 空间同构映射：利用 MDS 算法将该距离矩阵映射至连续的低维流形空间，生成各状态的坐标表征；
        \item 内生性约束：不同于常规的随机初始化，本方案直接将 MDS 生成的坐标映射为 Transformer 嵌入层（Embedding Layer）的初始权重，核心思路是“如果两个状态在物理上（或距离矩阵中）很接近，那么它们在 Embedding 空间中的初始位置也应该很接近。”
    \end{itemize}
    % 该路线通过在模型输入端建立有效的局部连续性约束，使 Transformer 能够原生感知状态间的邻近性，显著优化了搜索空间，从根本上抑制了非法轨迹的生成并提升了采样效率。
    % 这种方法实现了模型隐空间与物理拓扑空间的“同构”：在模型尚未接触训练数据之前，其 Embedding 向量间的欧氏距离已在数学上等价于物理空间距离。这使得 Transformer 的自注意力机制天然地获得了一种“局部性偏置（Inductive Bias）”，能够有效抑制跨越物理鸿沟的非法轨迹生成，显著提升了模型在复杂约束下的收敛效率与预测一致性。
\end{itemize}

\subsection*{2. 技术细节}
\begin{itemize}
        \item \textbf{基准版本 Vanilla}：使用标准Transformer架构，嵌入层随机初始化；
        \item \textbf{冻结版本 W/ IEF}：嵌入层保持空间拓扑不变，模型仅学习时序转移概率；
        \item \textbf{可训练版本 W/ IET}：以状态距离矩阵为强先验，通过数据驱动微调状态在隐空间的位置。
    \end{itemize}

\begin{figure}[ht]
    \centering
    \begin{minipage}{0.71\linewidth}
        \centering
        \includegraphics[width=\linewidth]{../../plot/training_comparison.png}
        Comparison of Training Loss and Total Mean Error
    \end{minipage}
    \begin{minipage}{0.265\linewidth}
        \centering
        \includegraphics[width=\linewidth]{../../plot/prediction_results.png}
        Ablation Experiments
    \end{minipage}
    \caption{Training curves and ablation experiments with different mechanisms}
    \label{fig:01}
\end{figure}

\newpage
图\ref{fig:01}展示了基准模型与两种不同策略下嵌入层初始化的训练曲线以及模型在测试集表现的对比。图\ref{fig:01}左侧两子图代表训练过程的Loss曲线与平均误差，横轴代表训练次数Epoch=20，相较于基准模型，使用状态距离矩阵进行Embedding初始化对训练效率与预测性能有一定提升，且嵌入层权重可训练版本模型性能更好；右图图中横轴代表测试集样本id（经预测误差排序），所示的消融实验中使用了嵌入层初始化机制的模型在大部分样本上具有更低的预测误差，验证了所设计机制的有效性。

\begin{figure}[ht]
    \centering
    \begin{minipage}{0.3\linewidth}
        \centering
        \includegraphics[width=\linewidth]{../../output/figures/pred_FL/traj_analysis_10.pdf}
        Vanilla
    \end{minipage}
    \hfill
    \begin{minipage}{0.3\linewidth}
        \centering
        \includegraphics[width=\linewidth]{../../output/figures/pred_FL_embedF/traj_analysis_10.pdf}
        W/ IEF
    \end{minipage}
    \hfill
    \begin{minipage}{0.3\linewidth}
        \centering
        \includegraphics[width=\linewidth]{../../output/figures/pred_FL_embedT/traj_analysis_10.pdf}
        W/ IET
    \end{minipage}
    \caption{test id: 10}
    \label{fig:02}
\end{figure}

图\ref{fig:02}-图\ref{fig:08}展示了不同测试案例下的预测结果（在状态值地形上的分布），根据测试样本的适应值在颜色条所对应的着色序列代表输入序列（决策过程中的历史序列），绿色虚线序列代表测试样本的实际后续状态序列，紫色序列代表模型预测结果（未来预测序列）。总体来看，冻结版本 W/ IEF与可训练版本 W/ IET的嵌入层初始化对预测结果有一定提升。一项符合表征学习预期的现象：冻结版本的嵌入为模型确立了稳健的嵌入空间基准，保证了其预测的逻辑下限；而可训练版本的的嵌入通过赋予模型微调表征的能力，进一步提升了预测精度。综合考虑两种策略的嵌入层初始化，这种通过问题特定下的状态距离矩阵构建嵌入层对训练效率与准确度的提升具有明显帮助。通过该技术构建一种未来序列预测模型，对辅助后续决策过程具有研究潜力。

目前对所提机制进行了初步尝试，后续考虑针对以下方向进行改进：
\begin{enumerate}
  \item 如何进一步提升预测模型的准确性（机制改良、模型超参调整）；
  \item 目前只针对状态序列进行预测，如若预测准确率高，也可以通过反向查表或统计等方式进行未来动作推荐，以直接辅助决策；是否可以将state-action的因果关系直接通过模型学习方式进行预测，以达到端到端的态势预测与动作推荐双层任务。
\end{enumerate}

\begin{figure}[ht]
    \centering
    \begin{minipage}{0.3\linewidth}
        \centering
        \includegraphics[width=\linewidth]{../../output/figures/pred_FL/traj_analysis_8.pdf}
        Vanilla
    \end{minipage}
    \hfill
    \begin{minipage}{0.3\linewidth}
        \centering
        \includegraphics[width=\linewidth]{../../output/figures/pred_FL_embedF/traj_analysis_8.pdf}
        W/ IEF
    \end{minipage}
    \hfill
    \begin{minipage}{0.3\linewidth}
        \centering
        \includegraphics[width=\linewidth]{../../output/figures/pred_FL_embedT/traj_analysis_8.pdf}
        W/ IET
    \end{minipage}
    \caption{test id: 8}
    \label{fig:03}
\end{figure}

\begin{figure}[ht]
    \centering
    \begin{minipage}{0.3\linewidth}
        \centering
        \includegraphics[width=\linewidth]{../../output/figures/pred_FL/traj_analysis_9.pdf}
        Vanilla
    \end{minipage}
    \hfill
    \begin{minipage}{0.3\linewidth}
        \centering
        \includegraphics[width=\linewidth]{../../output/figures/pred_FL_embedF/traj_analysis_9.pdf}
        W/ IEF
    \end{minipage}
    \hfill
    \begin{minipage}{0.3\linewidth}
        \centering
        \includegraphics[width=\linewidth]{../../output/figures/pred_FL_embedT/traj_analysis_9.pdf}
        W/ IET
    \end{minipage}
    \caption{test id: 9}
    \label{fig:04}
\end{figure}

\begin{figure}[ht]
    \centering
    \begin{minipage}{0.3\linewidth}
        \centering
        \includegraphics[width=\linewidth]{../../output/figures/pred_FL/traj_analysis_11.pdf}
        Vanilla
    \end{minipage}
    \hfill
    \begin{minipage}{0.3\linewidth}
        \centering
        \includegraphics[width=\linewidth]{../../output/figures/pred_FL_embedF/traj_analysis_11.pdf}
        W/ IEF
    \end{minipage}
    \hfill
    \begin{minipage}{0.3\linewidth}
        \centering
        \includegraphics[width=\linewidth]{../../output/figures/pred_FL_embedT/traj_analysis_11.pdf}
        W/ IET
    \end{minipage}
    \caption{test id: 11}
    \label{fig:05}
\end{figure}

\begin{figure}[ht]
    \centering
    \begin{minipage}{0.3\linewidth}
        \centering
        \includegraphics[width=\linewidth]{../../output/figures/pred_FL/traj_analysis_13.pdf}
        Vanilla
    \end{minipage}
    \hfill
    \begin{minipage}{0.3\linewidth}
        \centering
        \includegraphics[width=\linewidth]{../../output/figures/pred_FL_embedF/traj_analysis_13.pdf}
        W/ IEF
    \end{minipage}
    \hfill
    \begin{minipage}{0.3\linewidth}
        \centering
        \includegraphics[width=\linewidth]{../../output/figures/pred_FL_embedT/traj_analysis_13.pdf}
        W/ IET
    \end{minipage}
    \caption{test id: 13}
    \label{fig:06}
\end{figure}

\begin{figure}[ht]
    \centering
    \begin{minipage}{0.3\linewidth}
        \centering
        \includegraphics[width=\linewidth]{../../output/figures/pred_FL/traj_analysis_26.pdf}
        Vanilla
    \end{minipage}
    \hfill
    \begin{minipage}{0.3\linewidth}
        \centering
        \includegraphics[width=\linewidth]{../../output/figures/pred_FL_embedF/traj_analysis_26.pdf}
        W/ IEF
    \end{minipage}
    \hfill
    \begin{minipage}{0.3\linewidth}
        \centering
        \includegraphics[width=\linewidth]{../../output/figures/pred_FL_embedT/traj_analysis_26.pdf}
        W/ IET
    \end{minipage}
    \caption{test id: 26}
    \label{fig:07}
\end{figure}

\begin{figure}[ht]
    \centering
    \begin{minipage}{0.3\linewidth}
        \centering
        \includegraphics[width=\linewidth]{../../output/figures/pred_FL/traj_analysis_20.pdf}
        Vanilla
    \end{minipage}
    \hfill
    \begin{minipage}{0.3\linewidth}
        \centering
        \includegraphics[width=\linewidth]{../../output/figures/pred_FL_embedF/traj_analysis_20.pdf}
        W/ IEF
    \end{minipage}
    \hfill
    \begin{minipage}{0.3\linewidth}
        \centering
        \includegraphics[width=\linewidth]{../../output/figures/pred_FL_embedT/traj_analysis_20.pdf}
        W/ IET
    \end{minipage}
    \caption{test id: 20}
    \label{fig:08}
\end{figure}

% \centerline{\rule{\linewidth}{1pt}}
% \section*{\textbf{未来两周的研究计划 \\ {\fontsize{12pt}{4pt}\selectfont (Research work to do in the next two weeks)}}}
% \begin{enumerate}
%     \item 
% \end{enumerate}





\end{document}


